%% ============================================================
%% REFERÊNCIAS BIBLIOGRÁFICAS
%% ============================================================

\chapter{Referências Bibliográficas}

\begin{thebibliography}{99}

%% --- Documentos Normativos ---

\bibitem{bncc}
BRASIL. Ministério da Educação. \textbf{Base Nacional Comum Curricular (BNCC)}.
Brasília: MEC, 2018. Disponível em: \url{http://basenacionalcomum.mec.gov.br/}

\bibitem{dcnc}
BRASIL. Resolução CNE/CEB nº 2, de 22 de dezembro de 2012.
\textbf{Diretrizes Curriculares Nacionais do Ensino Fundamental de 9 (nove) anos}.
Diário Oficial da União, Brasília, 24 dez. 2012. Seção 1, p.~37.

\bibitem{ldb}
BRASIL. Lei nº 9.394, de 20 de dezembro de 1996.
\textbf{Estabelece as diretrizes e bases da educação nacional}.
Diário Oficial da União, Brasília, 23 dez. 1996. Seção 1, p.~27.833.

\bibitem{mecee}
BRASIL. Ministério da Educação. \textbf{Parecer CNE/CEB nº 7/2010.
Orientações Curriculares para o Ensino Médio}. Brasília: MEC, 2010.

%% --- Avaliação e Avaliações Padronizadas ---

\bibitem{caed}
CAED/UFJF. \textbf{Matriz de Referência para Avaliação ---
Língua Portuguesa}. Juiz de Fora: CAEd, 2019. Disponível em:
\url{https://www.saeb.inep.gov.br/}

\bibitem{inep_saeb}
INEP — Instituto Nacional de Estudos e Pesquisas Educacionais Anísio Teixeira.
\textbf{SAEB — Sistema de Avaliação da Educação Básica: Manual do Avaliador}.
Brasília: INEP, 2019. Disponível em: \url{https://www.gov.br/inep/pt-br}

\bibitem{inep_provabrasil}
INEP — Instituto Nacional de Estudos e Pesquisas Educacionais Anísio Teixeira.
\textbf{Prova Brasil: Manual do Aplicador}. Brasília: INEP, 2021.
Disponível em: \url{https://www.gov.br/inep/pt-br}

\bibitem{spaece}
FUNDAÇÃO CEARENSE DE APOIO AO DESENVOLVIMENTO CIENTÍFICO (FUNCAP).
\textbf{SAEB — Padrões de Desempenho em Língua Portuguesa}.
Fortaleza: FUNCAP, 2019.

%% --- Métodos e Abordagens Pedagógicas ---

\bibitem{kumon}
KUMON Institute of Education. \textbf{The Kumon Method}.
Tóquio: Kumon, 2020. Disponível em: \url{https://www.kumon.com/}

\bibitem{ehri}
EHRI, L. C. Learning to Read Words: Theory, Findings, and Issues.
\textit{Scientific Studies of Reading}, v.~9, n.~2, p.~167--188, 2005.

\bibitem{ferreiro}
FERREIRO, E.; TEBEROSKY, A. \textbf{Psicogênese da Língua Escrita}.
Porto Alegre: Artes Médicas, 1986.

\bibitem{morais}
MORAIS, A. \textbf{Alfabetização: Uma Abordagem Baseada no Processo
de Aquisição da Escrita}. São Paulo: Ática, 2007.

\bibitem{soares}
SOARES, M. \textbf{Alfabetização e Letramento}. São Paulo: Ática, 2003.

\bibitem{moretto}
MORETTO, V. \textbf{Letramento e Leitura}. São Paulo: Ática, 2004.

\bibitem{candeloro}
CANDELORO, L. \textbf{Alfabetização na Prática}. São Paulo: Cortez, 2003.

\bibitem{mortatti}
MORTATTI, H. \textbf{Ofício de Aprender a Ler e Escrever}.
São Paulo: Claraluz, 2004.

\bibitem{emilia}
EMÍLIA. \textbf{Da Cor ao Letramento}. São Paulo: Moderna, 1998.

\bibitem{pontes}
PONTES, C. A. \textbf{Alfabetização no Brasil: Histórias, Políticas
e Práticas Pedagógicas}. São Paulo: Ática, 2010.

\bibitem{gallindo}
GALLINDO, C. C. T. \textbf{Alfabetização no Brasil: O Estado da Arte}.
Campinas: Mercado das Letras, 2004.

\bibitem{motta}
MOTTA, E. S. \textbf{Alfabetização: Histórias e Narrativas}.
Campinas: Mercado das Letras, 2003.

%% --- Referências Teóricas ---

\bibitem{vygotsky}
VYGOTSKY, L. S. \textbf{Pensamento e Linguagem}. São Paulo: Martins
Fontes, 1989.

\bibitem{piaget}
PIAGET, J. \textbf{A Linguagem e o Pensamento da Criança}.
São Paulo: Mouro Brigido, 2001.

\bibitem{freire}
FREIRE, P. \textbf{Pedagogia do Oprimido}. Rio de Janeiro: Paz e Terra,
1987.

\bibitem{paivio}
PAIVIO, A. \textbf{Mental Representations: A Dual Coding Approach}.
Oxford: Oxford University Press, 1986.

\bibitem{kleiman}
KLEIMAN, A. B. \textbf{Textos e Leituras}. Campinas: Pontes, 2004.

\bibitem{koch}
KOCH, I. G. V. \textbf{Introdução à Linguística do Texto}.
São Paulo: Ática, 2004.

\bibitem{marques}
MARQUES, I. \textbf{Leitura e Escrita no Cotidiano Escolar}.
São Paulo: Moderna, 2004.

\end{thebibliography}

%% ============================================================
%% APÊNDICES
%% ============================================================

\appendix

%% ============================================================
%% APÊNDICE A — TABELA DE CORRESPONDÊNCIA FONEMA-GRAFEMA
%% ============================================================

\chapter{Tabela de Correspondência Fonema-Grafema Completa}

A tabela a seguir apresenta as correspondências entre as letras do
alfabeto português e seus fonemas. Para cada letra, indicam-se as
variantes fonéticas e exemplos de palavras em que cada correspon­dência
ocorre. Esta tabela é referência para a elaboração de atividades de
consciência fonológica e de reconhecimento grafema--fonema.

\begin{table}[h]
\centering
\caption{Tabela de Correspondência Fonema-Grafema --- Alfabeto Português}
\label{tab:fonema-grafema}
\renewcommand{\arraystretch}{1.25}
\begin{tabular}{|c|c|p{3.5cm}|p{3.5cm}|}
\hline
\textbf{Letra} & \textbf{Fonema(s)} & \textbf{Exemplos} & \textbf{Posição} \\
\hline
A & /a/ & água, avião, bola & V\'{a}ria \\
\hline
B & /b/ & bola, boi, umbrella & In\'{i}cio e meio \\
\hline
C & /k/ (a, o, u) & casa, cesto, coco & V\'{a}ria \\
  & /s/ (e, i) & cesto, cima, cedo & \\
\hline
Ç & /s/ & preço, força, laço & Meio e final \\
\hline
D & /d/ & dente, dado, dados & In\'{i}cio e meio \\
\hline
E & /e/ (aberto) & elefante, escola & V\'{a}ria \\
  & /e/ (fechado) & pé, café & \\
\hline
F & /f/ & fogo, fada, café & In\'{i}cio, meio e final \\
\hline
G & /g/ (a, o, u) & gato, goiaba, gula & V\'{a}ria \\
  & /zh/ (e, i) & gelo, girafa & \\
\hline
H & --- (mudo) & hotel, homem & In\'{i}cio (ap\'{o}s C e L);
n\~{a}o tem som pr\'{o}prio \\
\hline
I & /i/ & igreja, ipê, saí & V\'{a}ria \\
\hline
J & /zh/ & jabuti, janela, biju & In\'{i}cio e final \\
\hline
K & /k/ & kiwi, whisky & Uso restrito (palavras estrangeiras) \\
\hline
L & /l/ & lua, limão, sal & In\'{i}cio e final \\
  & /u/ (mudo) & sal, mal, papel & \\
\hline
M & /m/ & macaco, mesa, jogo & In\'{i}cio e meio \\
\hline
N & /n/ & navio, nota, ANO & In\'{i}cio e final \\
\hline
O & /o/ (aberto) & ovo, olho, bolo & V\'{a}ria \\
  & /o/ (fechado) & m\'{o}vel, p\'{o}r & \\
\hline
P & /p/ & pato, pedra, sapato & In\'{i}cio e final \\
\hline
Q & /k/ & queijo, quatro, quota & In\'{i}cio (antes de U) \\
\hline
R & /R/ (gutural) & rato, ramo, roupa & In\'{i}cio (produzido na
garganta) \\
  & /r/ (alveolar) & cara, muro, flor & Meio e final \\
\hline
S & /s/ (alveolar) & sapo, sopa, masa & In\'{i}cio e final \\
  & /z/ (alveolar) & casa, mesa, raso & Meio (entre vogais) \\
\hline
T & /t/ & tigre, teto, mato & In\'{i}cio e final \\
\hline
U & /u/ & uva, urso, sal & V\'{a}ria \\
\hline
V & /v/ & vaso, vento, chave & In\'{i}cio e final \\
\hline
W & /w/ & whisky, web & Uso restrito (palavras estrangeiras) \\
\hline
X & /sh/ & xícara, xarope, Choque & V\'{a}ria \\
  & /ks/ & texto, fixo & \\
  & /z/ & exame, exército & \\
\hline
Y & /i/ & kayak, Boeing & Uso restrito (palavras estrangeiras) \\
\hline
Z & /z/ & zebra, zero, casa & In\'{i}cio e final \\
\hline
\end{tabular}
\end{table}

%% ============================================================
%% APÊNDICE B — SEQUÊNCIA DE LETRAS POR BLOCO
%% ============================================================

\chapter{Sequência de Letras por Bloco de Apresentação}

A ordem de apresentação das letras segue o m\'{e}todo fon\'{e}tico,
come\c{c}ando pelas letras mais frequentes e de maior facilidade
de articula\c{c}\~{a}o. Cada bloco \'{e} introduzido em um volume
espec\'{i}fico do m\'{e}tdo h\'{i}brido Fonico-Kumon.

\begin{table}[h]
\centering
\caption{Sequ\^{e}ncia de Letras por Bloco de Apresenta\c{c}\~{a}o}
\label{tab:sequencia-letras}
\renewcommand{\arraystretch}{1.35}
\begin{tabular}{|c|c|p{5cm}|c|}
\hline
\textbf{Bloco} & \textbf{Letras} & \textbf{Justificativa} & \textbf{Volume} \\
\hline
\rowcolor{warning!15}
1 &
a, o, e, i, u, m, s, t, r, l &
Vogais (5) + consoantes de maior frequ\^{e}ncia e facilidade
de articula\c{c}\~{a}o (5). Formam s\'{i}labas diretas (CV). &
Vol.~1 \\
\hline
\rowcolor{accent!10}
2 &
c, d, p, n, b, g, f, v &
Consoantes do segundo n\'{i}vel de frequ\^{e}ncia. Permitem
a forma\c{c}\~{a}o de novas s\'{i}labas diretas e o in\'{i}cio
de s\'{i}labas indiretas (VC). &
Vol.~2 \\
\hline
\rowcolor{primary!10}
3 &
j, q, x, z, h, y, w, k &
Letras de menor frequ\^{e}ncia e uso mais restrito. Q sempre
aparece com U. H \'{e} mudo. K, W e Y s\~{a}o usadas em
palavras estrangeiras. &
Vol.~3 \\
\hline
\rowcolor{success!10}
4 &
ch, lh, nh, rr, ss, ll, gh &
Digrafos: duas letras que representam um \'{u}nico fonema.
Consoantes duplas: RR e SS (o\={a}s representam um fonema
diferente da letra simples). &
Vol.~3 \\
\hline
5 &
ai, ei, oi, ui, au, eu, ou, ia, ie, io, iu, ua, ue, uo &
Ditongos: encontros voc\'{a}licos em uma mesma s\'{i}laba.
Divididos em abertos (terminados em I ou U) e fechados
(terminados em L ou R). &
Vol.~4 \\
\hline
\end{tabular}
\end{table}

%% ============================================================
%% APÊNDICE C — GLOSSÁRIO DE TERMOS TÉCNICOS
%% ============================================================

\chapter{Glossário de Termos Técnicos}

O glossário a seguir re\'{u}ne os principais termos t\'{e}cnicos
utilizados neste livro did\'{a}tico, incluindo termos das Diretrizes
Curriculares Nacionais do Ensino Fundamental (DCNC) e da Base
Nacional Comum Curricular (BNCC).

\begin{description}[leftmargin=2.5cm, style=nextline]

  \item[Alfabetização] Processo de aprendizagem da leitura e da
  escrita. Envolve a aquisi\c{c}\~{a}o do princ\'{i}pio alfab\'{e}tico
  e o desenvolvimento da consci\^{e}ncia fonol\'{o}gica.

  \item[Alfabético (n\'{i}vel)] N\'{i}vel de alfabetiza\c{c}\~{a}o
  em que o aluno domina o princ\'{i}pio alfab\'{e}tico: relaciona
  fonemas a grafemas em todas as posi\c{c}\~{o}es e l\'{e} com flu\^{e}ncia.

  \item[BNCC] Base Nacional Comum Curricular --- documento que
  estabelece os objetivos de aprendizagem para a educa\c{c}\~{a}o
  b\'{a}sica brasileira (BRASIL, 2018).

  \item[CAEd] Centro de Avaliação Educacional --- institui\c{c}\~{a}o
  respons\'{a}vel pelas avalia\c{c}\~{o}es da Prova Brasil e do SAEB.

  \item[Consci\^{e}ncia Fonol\'{o}gica] Habilidade de perceber e
  manipular os sons da l\'{i}ngua oral: rima, alitera\c{c}\~{a}o,
  segmenta\c{c}\~{a}o sil\'{a}bica e fonêmica.

  \item[DCNC] Diretrizes Curriculares Nacionais do Ensino Fundamental
  --- resolu\c{c}\~{a}o CNE/CEB nº 2/2012 que estabelece princ\'{i}pios,
  fundamentos e procedimentos para o ensino fundamental de 9 anos.

  \item[Decodifica\c{c}\~{a}o] Processo de converter letras (grafemas)
  em sons (fonemas) para ler palavras.

  \item[Descritor] Habilidade espec\'{i}fica avaliada em itens de
  prova. Na matriz CAEd, cada descritor corresponde a uma habilidade
  de L\'{i}ngua Portuguesa (D1, D2, D3, D4, D5).

  \item[Diferencia\c{c}\~{a}o] Estrat\'{e}gias pedag\'{o}gicas que
  permitem atender a diferentes n\'{i}veis de habilidade dentro da
  mesma turma.

  \item[Ditongo] Encontro de duas vogais na mesma s\'{i}laba (ex.: AI
  em ``cair'', OU em ``couve'').

  \item[Fonema] Unidade m\'{i}nima de som da fala que distingue palavras
  (ex.: /m/ em ``mas'' e /n/ em ``nas'').

  \item[Grafema] Representa\c{c}\~{a}o gr\'{a}fica de um fonema
  (letra ou grupo de letras).

  \item[Hip\'{o}tese Alfab\'{e}tica] Compreens\'{e}ncia de que cada
  letra corresponde a um som espec\'{i}fico e de que a combina\c{c}\~{a}o
  de letras permite ler qualquer palavra.

  \item[Hip\'{o}tese Sil\'{a}bica] Compreens\'{e}ncia de que as letras
  representam s\'{i}labas e n\~{a}o fonemas isolados.

  \item[Letramento] Capacidade de usar a leitura e a escrita na vida
  cotidiana, em pr\'{a}ticas sociais significativas.

  \item[LDB] Lei de Diretrizes e Bases da Educa\c{c}\~{a}o Nacional
  (Lei nº 9.394/1996) --- diploma que organiza o sistema educacional
  brasileiro.

  \item[Mastery-based] Aprendizagem onde o aluno s\'{o} avan\c{c}a
  quando domina o conte\'{u}do atual (princ\'{i}pio do Kumon).

  \item[Pré-sil\'{a}bico (n\'{i}vel)] N\'{i}vel de alfabetiza\c{c}\~{a}o
  em que o aluno ainda n\~{a}o compreende que as letras representam sons.
  Reconhece letras, mas n\~{a}o estabelece rela\c{c}\~{a}o
  fonema--grafema.

  \item[Progress\~{a}o Granular] Decomposi\c{c}\~{a}o de habilidades
  em micro-passos incrementais, permitindo que o aluno avance
  gradualmente.

  \item[Prova Brasil] Avalia\c{c}\~{a}o em larga escala do INEP que
  mede o desempenho de alunos do 5º e 9º anos do ensino fundamental
  em L\'{i}ngua Portuguesa e Matem\'{a}tica.

  \item[SAEB] Sistema de Avalia\c{c}\~{a}o da Educa\c{c}\~{a}o B\'{a}sica
  --- avalia\c{c}\~{a}o em larga escala realizada a cada quatro anos
  pelo INEP.

  \item[Sílaba Composta] S\'{i}laba com dois ou mais fonemas em uma
  mesma posi\c{c}\~{a}o (ex.: TRA, PRA, BLA).

  \item[Sílaba Direta] S\'{i}laba formada por consoante + vogal
  (ex.: MA, CE, TI).

  \item[Sílaba Indireta] S\'{i}laba formada por vogal + consoante
  (ex.: AM, ES, ON).

  \item[Sil\'{a}bico (n\'{i}vel)] N\'{i}vel de alfabetiza\c{c}\~{a}o
  em que o aluno compreende que letras representam sons, mas ainda
  n\~{a}o domina a rela\c{c}\~{a}o exata entre fonema e grafema em
  todas as posi\c{c}\~{o}es.

  \item[Sil\'{a}bico-Alfab\'{e}tico (n\'{i}vel)] N\'{i}vel de
  alfabetiza\c{c}\~{a}o em que o aluno utiliza estrat\'{e}gias mistas:
  decodifica parcialmente e utiliza pistas do contexto.

  \item[SPAECE] Sistema Permanente de Avalia\c{c}\~{a}o da Educa\c{c}\~{a}o
  B\'{a}sica do Ceará.

\end{description}

%% ============================================================
%% APÊNDICE D — MODELO DE ITEM CAEd
%% ============================================================

\chapter{Modelos de Itens CAEd}

Os itens abaixo seguem o formato da Matriz de Referência do SAEB/Prova Brasil,
elaborados pelo CAEd/UFJF. Cada item \'{e} associado a um descritor
espec\'{i}fico e avalia uma habilidade de L\'{i}ngua Portuguesa.

%% ----------------------------------------------------------------
\section{Item CAEd --- D1: Identifica\c{c}\~{a}o de Fonemas}
%% ----------------------------------------------------------------

\begin{exercicio}[Item CAEd --- Descritor D1: Identifica\c{c}\~{a}o de Fonemas]

\textbf{Enunciado:} Na palavra \textbf{``mesa}'', qual \'{e} o som
que a letra M representa?

\begin{enumerate}[label=\alph*)]
  \item /a/ (como em ``\textbf{a}guia'')
  \item /m/ (como em ``\textbf{m}acaco'')
  \item /s/ (como em ``\textbf{s}apo'')
  \item /z/ (como em ``me\textbf{z}a'')
\end{enumerate}

\textbf{Gabarito:} \textbf{b)}

\textbf{Coment\'{a}rio pedag\'{o}gico:} Este item avalia se o aluno
consegue identificar o fonema associado a uma letra espec\'{i}fica
dentro de uma palavra. O aluno deve reconhecer que M = /m/ em todas
as posi\c{c}\~{o}es da palavra.
\end{exercicio}

%% ----------------------------------------------------------------
\section{Item CAEd --- D2: Leitura de Palavras}
%% ----------------------------------------------------------------

\begin{exercicio}[Item CAEd --- Descritor D2: Leitura de Palavras]

\textbf{Enunciado:} Leia a palavra destacada e escolha a alternativa
que indica corretamente o que ela significa:

\begin{center}
\textbf{A menina \underline{corre} no parque.}
\end{center}

\begin{enumerate}[label=\alph*)]
  \item Dorme
  \item Come
  \item Corre
  \item Pula
\end{enumerate}

\textbf{Gabarito:} \textbf{c)}

\textbf{Coment\'{a}rio pedag\'{o}gico:} Este item avalia a
decodifica\c{c}\~{a}o de palavras simples com s\'{i}labas diretas
(CV). O aluno deve decodificar a palavra ``corre'' e associar ao
seu significado.
\end{exercicio}

%% ----------------------------------------------------------------
\section{Item CAEd --- D3: Leitura de Frases}
%% ----------------------------------------------------------------

\begin{exercicio}[Item CAEd --- Descritor D3: Leitura de Frases]

\textbf{Enunciado:} Leia a frase e assinale o que est\'{a} escrito:

\begin{center}
\textbf{O menino est\'{a} brincando no quintal.}
\end{center}

\begin{enumerate}[label=\alph*)]
  \item O menino est\'{a} estudando na escola.
  \item O menino est\'{a} brincando no quintal.
  \item A menina est\'{a} brincando no quintal.
  \item O menino est\'{a} dormindo no quintal.
\end{enumerate}

\textbf{Gabarito:} \textbf{b)}

\textbf{Coment\'{a}rio pedag\'{o}gico:} Este item avalia a
compreens\~{a}o de frases simples. O aluno deve ler a frase
completamente e identificar a que corresponde ao texto apresentado,
diferenciando sujeito (menino/menina) e a\c{c}\~{a}o
(brincando/estudando/dormindo).
\end{exercicio}

%% ----------------------------------------------------------------
\section{Item CAEd --- D4: Separa\c{c}\~{a}o Sil\'{a}bica}
%% ----------------------------------------------------------------

\begin{exercicio}[Item CAEd --- Descritor D4: Separa\c{c}\~{a}o Sil\'{a}bica]

\textbf{Enunciado:} Assinale a alternativa em que a palavra
\textbf{``sala}'' est\'{a} separada corretamente em s\'{i}labas:

\begin{enumerate}[label=\alph*)]
  \item sa--la
  \item sal--a
  \item s--ala
  \item sa--l--a
\end{enumerate}

\textbf{Gabarito:} \textbf{a)}

\textbf{Coment\'{a}rio pedag\'{o}gico:} Este item avalia se o aluno
identifica a fronteira sil\'{a}bica em palavras com s\'{i}labas
diretas (CV). A palavra ``sala'' \'{e} dividida em SA + LA.
\end{exercicio}

%% ----------------------------------------------------------------
\section{Item CAEd --- D5: Produ\c{c}\~{a}o de Frases}
%% ----------------------------------------------------------------

\begin{exercicio}[Item CAEd --- Descritor D5: Produ\c{c}\~{a}o de Frases]

\textbf{Enunciado:} Assinale a frase que est\'{a} escrita corretamente:

\begin{enumerate}[label=\alph*)]
  \item O gato mia.
  \item O gato mia
  \item o gato mia.
  \item O gato mia .
\end{enumerate}

\textbf{Gabarito:} \textbf{a)}

\textbf{Coment\'{a}rio pedag\'{o}gico:} Este item avalia se o aluno
reconhece as conven\c{c}\~{o}es b\'{a}sicas de escrita: in\'{i}cio
com letra mai\'{u}scula e ponto final. A op\c{c}\~{a}o (a) \'{e} a
\'{u}nica que atende a todas as conven\c{c}\~{o}es.
\end{exercicio}
